% Generated draft: requires mathematical and source review.
\documentclass[11pt]{article}
\usepackage[margin=25mm]{geometry}
\usepackage{fontspec}
\setmainfont{Latin Modern Roman}
\usepackage{amsmath,amssymb,mathtools,unicode-math}
\setmathfont{Latin Modern Math}
\usepackage{enumitem,needspace,xurl,hyperref}
\setlength{\parindent}{0pt}
\setlength{\parskip}{5pt}
\setlength{\emergencystretch}{3em}
\newcommand{\N}{\mathbb{N}}
\newcommand{\Z}{\mathbb{Z}}
\newcommand{\Q}{\mathbb{Q}}
\newcommand{\R}{\mathbb{R}}
\newcommand{\C}{\mathbb{C}}
\newcommand{\F}{\mathbb{F}}
\newcommand{\A}{\mathbb{A}}
\newcommand{\PP}{\mathbb{P}}
\newcommand{\E}{\mathbb{E}}
\newcommand{\eps}{\varepsilon}
\newcommand{\dd}{\,\mathrm{d}}
\begin{document}
\section*{UnsolvedMath ID: 701645}
\textbf{Source problem code:} catalogue entry
\subsection{Definitions and mathematical statement}
For a group \(H\), define its derived series by \(H^{(0)}=H\) and \(H^{(k+1)}=[H^{(k)},H^{(k)}]\). The group has derived length \(n\) if \(H^{(n)}=1\) and \(H^{(n-1)}\neq 1\). Let \(L_r\) be the free group on \(r\) generators. The free soluble group of rank \(r\) and derived length \(n\) is
\[
F=L_r/L_r^{(n)}.
\]
A group is \(m\)-generated if it has a generating set of at most \(m\) elements; minimality of this generating set is not assumed. Write \(\operatorname{Th}(H)\) for the set of all first-order sentences true in \(H\), in the language of groups with multiplication, inversion, and identity.

For every pair of integers \(n\geq 3\) and \(r\geq 2\), every positive integer \(m\), and every \(m\)-generated group \(G\), does
\[
\operatorname{Th}(G)=\operatorname{Th}(F)
\quad\Longrightarrow\quad
G\cong F
\]
hold?

The supplied record states that the answer is affirmative when \(m\leq r\). It also states that the analogous question has an affirmative answer for \(n=2\), a case outside the main question's stated range \(n\geq 3\). These affirmative-case assertions and the record's open-status label are supplied source qualifications, not independently verified conclusions.
\subsection{Short English statement}
Must every finitely generated group with the same first-order theory as a finite-rank free soluble group of derived length at least three be isomorphic to that free soluble group? The source reports affirmative answers when \(m\leq r\), and for the analogous derived-length-two case.
\subsection{Sources}
\begin{itemize}
\item Local UnsolvedMath registry, supplied record with permanent ID \(701645\), exact source code \texttt{catalogue entry}, title \emph{Kourovka Notebook Problem 19.97}. Upstream revision: \texttt{d371ee917ac9647fbb06cd6fc529deec6f3dda82}. The record labels the problem open and was updated on \texttt{2026-10-16T14:00:00Z}; this is supplied registry metadata, not independent verification of its mathematical status.
\item E. I. Timoshenko, problem 19.97, proposed in 2018, in \emph{The Kourovka Notebook (Unsolved Problems in Group Theory)}, 19th issue (2018). The supplied background also identifies the 21st issue edition (2026), \texttt{arXiv:1401.0300v48}. The collection metadata names E. I. Khukhro and V. D. Mazurov as editors.
\end{itemize}
\end{document}
